\documentclass[12pt]{amsart}
\usepackage{amsmath,amssymb,amsthm}
\usepackage[margin=1in]{geometry}
\usepackage{enumitem}
\usepackage{booktabs}
\usepackage{hyperref}

\newtheorem{theorem}{Theorem}[section]
\newtheorem{lemma}[theorem]{Lemma}
\newtheorem{corollary}[theorem]{Corollary}
\newtheorem{proposition}[theorem]{Proposition}
\theoremstyle{definition}
\newtheorem{definition}[theorem]{Definition}
\theoremstyle{remark}
\newtheorem{remark}[theorem]{Remark}

\newcommand{\CC}{\mathbb{C}}
\newcommand{\RR}{\mathbb{R}}
\newcommand{\HH}{\mathbb{H}}
\newcommand{\ZZ}{\mathbb{Z}}
\newcommand{\CP}{\mathbb{C}P}
\newcommand{\tr}{\operatorname{Tr}}
\newcommand{\agut}{\alpha_{\mathrm{GUT}}}

\title{From Frobenius to Gauge:\\
  $\CC$ as the Unique Substrate,\\
  $\CC^5 = \CC^2 \oplus \CC^3$ as the Unique Structure}
\author{Mingu Jeong}
\address{Independent Researcher}
\author{Mingu Jeong}
\address{Anthropic}
\date{\today}

\begin{document}
\begin{abstract}
Starting from the requirement that pairwise
relations $G_{ij}$ between $N$ points take
values in a normed division algebra
supporting magnitude (geometry),
continuous phase (gauge forces),
a well-defined determinant (gravity),
and integer winding numbers (charge quantization),
we prove that the unique such algebra is $\CC$.
Combined with the chiral atomic decomposition
theorem ($\CC^5 = \CC^2 \oplus \CC^3$ is unique),
this determines the gauge group
$\mathrm{SU}(3) \times \mathrm{SU}(2) \times \mathrm{U}(1)$
and the unified coupling constant
$\agut = 6 / (25\pi^2)$,
the latter derived by three independent paths:
simplex geometry, random matrix theory,
and number-theoretic coprimality.
\end{abstract}

\maketitle
\tableofcontents

%=====================================================================
\section{Introduction}
%=====================================================================

The Standard Model of particle physics is built on
the gauge group
$\mathrm{SU}(3) \times \mathrm{SU}(2) \times \mathrm{U}(1)$,
the complex numbers $\CC$ as the field of quantum mechanics,
and a specific set of coupling constants
determined by experiment.
None of these inputs are derived within the Standard Model;
they are axioms.

We show that all three are theorems.
The argument proceeds in three stages:

\begin{enumerate}[label=(\roman*)]
\item \textbf{The substrate} (Section~\ref{sec:substrate}).
  Four necessary conditions for a universe
  with geometry, forces, gravity, and discrete
  charges uniquely select $\CC$
  from the Frobenius classification of
  division algebras.

\item \textbf{The structure}
  (Section~\ref{sec:structure}).
  The chiral atomic decomposition theorem
  \cite{Paper1} determines
  $\CC^5 = \CC^2 \oplus \CC^3$,
  yielding
  $\mathrm{SU}(3) \times \mathrm{SU}(2)
  \times \mathrm{U}(1)$.

\item \textbf{The coupling}
  (Sections~\ref{sec:coupling}--\ref{sec:three_paths}).
  The unified coupling constant
  $\agut = 6/(25\pi^2)$ is derived
  by three independent methods,
  all originating from $\CC$.
\end{enumerate}

%=====================================================================
\section{The Unique Substrate: $\CC$}
\label{sec:substrate}
%=====================================================================

\subsection{The axiom}

We begin with a minimal assumption:

\begin{quote}
\textbf{Axiom.}
There exist $N$ points.
Between every ordered pair $(i, j)$
there exists a relation $G_{ij}$
taking values in an algebra $\mathbb{K}$.
\end{quote}

\subsection{Four necessary conditions}

For the relation matrix $G = (G_{ij})$
to encode geometry, forces, gravity,
and discrete charges,
$\mathbb{K}$ must satisfy:

\begin{enumerate}[label=\textbf{R\arabic*.}]
\item \textbf{Magnitude} (geometry):
  $|G_{ij}|$ is well-defined.
  Requires $\mathbb{K}$ normed.
\item \textbf{Phase} (gauge forces):
  $\arg(G_{ij})$ is continuous.
  Requires unit-norm elements to form
  a connected Lie group.
\item \textbf{Determinant} (gravity):
  $\det(G)$ is a single real number.
  Requires $\mathbb{K}$ commutative.
\item \textbf{Winding} (charge quantization):
  holonomy admits discrete winding numbers.
  Requires $\pi_1(\text{phase group}) \cong \ZZ$.
\end{enumerate}

\subsection{Algebraic proof}

\begin{theorem}[Uniqueness of $\CC$]
\label{thm:C_unique}
$\CC$ is the unique finite-dimensional
associative division algebra over $\RR$
satisfying R1--R4.
\end{theorem}

\begin{proof}
By the Frobenius theorem \cite{Frobenius},
the candidates are
$\{\RR,\, \CC,\, \HH\}$.
(The octonions $\mathbb{O}$ are excluded by
non-associativity.)

\medskip\noindent
\textbf{$\HH$ violates R3.}
The quaternions are non-commutative:
$\mathbf{i}\mathbf{j} = \mathbf{k}
\ne -\mathbf{k} = \mathbf{j}\mathbf{i}$.
The determinant of a matrix over $\HH$
is not well-defined as a single real number.

\medskip\noindent
\textbf{$\RR$ violates R2.}
The unit-norm elements of $\RR$ are
$\{+1, -1\} \cong \ZZ_2$,
a discrete group.
No continuous phase, no gauge forces.

\medskip\noindent
\textbf{$\CC$ satisfies all four.}
$|z|$ is a norm (R1).
$\arg(z) \in (-\pi, \pi]$
and $\mathrm{U}(1) = S^1$ is a
connected Lie group (R2).
$\CC$ is commutative (R3).
$\pi_1(S^1) = \ZZ$ (R4).
\end{proof}

\subsection{Topological confirmation}

\begin{theorem}
$\CC$ is the unique normed division algebra
whose phase group has $\pi_1 = \ZZ$.
\end{theorem}

\begin{proof}
The phase groups are spheres:
$\RR \to S^0$,
$\CC \to S^1$,
$\HH \to S^3$,
$\mathbb{O} \to S^7$.
Their fundamental groups:
$\pi_1(S^0) = 0$ (disconnected),
$\pi_1(S^1) = \ZZ$,
$\pi_1(S^3) = 0$ (simply connected),
$\pi_1(S^7) = 0$.
Only $\CC$ has $\pi_1 = \ZZ$.
\end{proof}

%=====================================================================
\section{The Unique Structure: $\CC^2 \oplus \CC^3$}
\label{sec:structure}
%=====================================================================

Having established $\mathbb{K} = \CC$,
the relation takes the form
$G_{ij} = \langle \psi_i, \psi_j \rangle$
with $\psi_i \in \CC^d$, $\|\psi_i\| = 1$.
The dimension $d$ and its decomposition
are determined by the chiral atomic
decomposition theorem \cite{Paper1}:

\begin{theorem}[{\cite[Theorem~5.1]{Paper1}}]
The unique chiral atomic decomposition is
$\CC^5 = \CC^2 \oplus \CC^3$.
\end{theorem}

This decomposition determines the gauge group:

\begin{corollary}[Gauge group]
\label{cor:gauge}
The structure group of
$\CC^5 = \CC^2 \oplus \CC^3$ is
\begin{equation}
\mathrm{S}(\mathrm{U}(3) \times \mathrm{U}(2))
\cong
\mathrm{SU}(3) \times \mathrm{SU}(2)
\times \mathrm{U}(1).
\end{equation}
The $\mathrm{U}(1)$ factor is the relative
phase between the two summands,
with $\pi_1(\mathrm{U}(1)) = \ZZ$
providing charge quantization (R4).
The tracelessness condition
$3 y_A + 2 y_B = 0$ fixes the
hypercharge ratio $y_A / y_B = -2/3$.
\end{corollary}

%=====================================================================
\section{The Born Rule as Algebraic Necessity}
\label{sec:born}
%=====================================================================

\begin{theorem}
The unique polynomial map
$f \colon \CC \to \RR_{\ge 0}$ of minimal degree
satisfying:
(i)~real-valued,
(ii)~symmetric ($f(G_{ij}) = f(G_{ji})$),
(iii)~non-negative,
(iv)~faithful ($f(z) = 0 \Leftrightarrow z = 0$),
is $f(z) = z\bar{z} = |z|^2$.
\end{theorem}

\begin{proof}
Any polynomial in $z$ alone is complex-valued
(violates (i)).
Any real polynomial in $z$ and $\bar{z}$
that is non-negative must contain the factor
$z\bar{z}$;
the minimal such polynomial is $|z|^2$
(degree~$2$).
\end{proof}

The observable overlap is therefore
$W_{ij} = |G_{ij}|^2 / d$:
the Born rule $P = |\langle \psi | \psi' \rangle|^2$
is a property of $\CC$, not an axiom
of quantum mechanics.

%=====================================================================
\section{Coupling Constants from $\CC^2 \oplus \CC^3$}
\label{sec:coupling}
%=====================================================================

\subsection{Binet--Cauchy channel counting}

The third exterior power
$\Lambda^3(\CC^5)$ decomposes under
$\CC^5 = \CC^3 \oplus \CC^2$ as:
\begin{equation}
\Lambda^3(\CC^5)
= \bigoplus_{k=0}^{2}
  \Lambda^{3-k}(\CC^3) \otimes \Lambda^k(\CC^2).
\end{equation}
The term $k = 3$ vanishes since
$\Lambda^3(\CC^2) = 0$.

With the lattice speed $c = n_T / (d_T / d_S)
= 2$ (derived from the $(3,2)$ asymmetry),
the $c$-weighted channel count is:
\begin{align}
k = 0: &\quad
  \tbinom{3}{3}\tbinom{2}{0} \cdot c^0
  = 1
  & \text{(pure-$A$ sector)}, \\
k = 1: &\quad
  \tbinom{3}{2}\tbinom{2}{1} \cdot c^1
  = 12
  & \text{(mixed $AAB$)}, \\
k = 2: &\quad
  \tbinom{3}{1}\tbinom{2}{2} \cdot c^2
  = 12
  & \text{(mixed $ABB$)}.
\end{align}

\begin{theorem}[Channel sum]
\label{thm:25}
$1 + 12 + 12 = 25 = d^2$.
This identity holds if and only if $c = 2$.
\end{theorem}

\subsection{The Basel propagator}

Each channel propagates via the
solid-angle propagator $D(n) = 1/n^2$
(in $n_A = 3$ spatial dimensions).
The total coupling per channel over
range $N_{\mathrm{eff}}$ is:
\begin{equation}
S(N_{\mathrm{eff}})
= \sum_{n=1}^{N_{\mathrm{eff}}} \frac{1}{n^2}.
\end{equation}
For infinite range:
$S(\infty) = \zeta(2) = \pi^2/6$.

\subsection{Topological horizon}

The propagation range $N_{\mathrm{eff}}$
is derived from rank constraints:

\begin{itemize}
\item \textbf{$A$-sector ($k=0$):}
  $\binom{3}{3} = 1$ configuration.
  Rank exhausted in 1 hop.
  $N_{\mathrm{eff}} = 1$, $S(1) = 1$.

\item \textbf{$B$-sector ($k=2$):}
  $\dim \CC^2 = 2$.
  Rank exhausted in 2 hops.
  $N_{\mathrm{eff}} = 2$, $S(2) = 5/4$.

\item \textbf{Cross-sector ($k=1$):}
  $\mathrm{U}(1)$ phase between sectors.
  No rank exhaustion.
  $N_{\mathrm{eff}} = \infty$,
  $S(\infty) = \pi^2/6$.
\end{itemize}

\subsection{The coupling constants}

\begin{theorem}[Inverse couplings]
\label{thm:couplings}
\begin{align}
1/\alpha_3 &= 1 \times (3^2 - 1) \times 1
  = 8, \\
1/\alpha_2 &= 12 \times 2 \times 5/4
  = 30, \\
1/\alpha_1 &= 12 \times 3 \times \pi^2/6
  = 6\pi^2 \approx 59.22.
\end{align}
The gauge multiplicities follow a single rule:
pure sectors get the adjoint ($n^2 - 1$),
mixed sectors get the fundamental ($n$).
\end{theorem}

%=====================================================================
\section{Three Independent Paths to $\agut$}
\label{sec:three_paths}
%=====================================================================

At the unification scale,
all $d^2 = 25$ channels see
$N_{\mathrm{eff}} = \infty$.
The unified coupling is:
\begin{equation}
\frac{1}{\agut}
= d^2 \times \zeta(2)
= \frac{25\pi^2}{6}
\approx 41.12.
\end{equation}

This constant is derived by three independent paths:

\subsection{Path 1: Simplex geometry (Binet--Cauchy)}

As shown in Section~\ref{sec:coupling}:
$d^2 = 25$ channels, each with propagator
$S(\infty) = \zeta(2)$.
Hence $1/\agut = 25 \cdot \pi^2/6$.

\subsection{Path 2: Random matrix theory (GUE)}

The Gram matrix $G = \Psi\Psi^{\dagger}$
with $\Psi \in \CC^{N \times d}$ belongs
to the Gaussian Unitary Ensemble
($\beta = 2$, forced by $\mathbb{K} = \CC$).
The two-point cluster function of the
sine kernel is
$Y_2(r) = (\sin \pi r / \pi r)^2$.
At short distance,
$R_2(r) \approx 2\zeta(2) r^2$.
Partitioning over $d^2$ channels:
$\agut = 1/(d^2 \zeta(2)) = 6/(25\pi^2)$.

\subsection{Path 3: Number theory (coprimality)}

The probability that two random integers
are coprime is
$P(\gcd(a,b) = 1) = 1/\zeta(2) = 6/\pi^2$
(Euler--Ces\`aro).
Among $d^2 = 25$ channels:
$\agut = P(\text{coprime})/d^2 = 6/(25\pi^2)$.

\begin{theorem}[GUT coupling as mathematical constant]
\label{thm:agut}
$\agut = 6/(25\pi^2)$ is a theorem,
derived independently by geometry,
analysis, and number theory.
It is not a measurement.
\end{theorem}

%=====================================================================
\section{Conclusion}
%=====================================================================

Starting from ``$N$ points with pairwise
relations,'' we have derived in sequence:
\begin{enumerate}
\item $\CC$ as the unique substrate
  (Frobenius, R1--R4),
\item $\CC^5 = \CC^2 \oplus \CC^3$ as the
  unique chiral structure \cite{Paper1},
\item $\mathrm{SU}(3) \times \mathrm{SU}(2)
  \times \mathrm{U}(1)$ as the gauge group,
\item $\agut = 6/(25\pi^2)$ as the unified coupling,
\item $1/\alpha_3 = 8$,\;
  $1/\alpha_2 = 30$,\;
  $1/\alpha_1 = 6\pi^2$
  as the low-energy couplings.
\end{enumerate}

The derivation of fermion masses,
cosmological observables, and further
predictions is presented in the
companion paper \cite{Paper3}.

\begin{thebibliography}{99}

\bibitem{Paper1}
M.~Jeong and Claude,
``Uniqueness of Chiral Atomic Decomposition
on $\CC^d$,'' companion paper (2026).

\bibitem{Paper3}
M.~Jeong and Claude,
``Zero-Parameter Predictions from
Simplex Geometry,'' companion paper (2026).

\bibitem{Frobenius}
F.~G.~Frobenius,
``\"Uber lineare Substitutionen und
bilineare Formen,''
J.~Reine Angew.~Math. \textbf{84},
1--63 (1877).

\bibitem{Hurwitz}
A.~Hurwitz,
``\"Uber die Composition der quadratischen
Formen,''
Nachr.~Ges.~Wiss.~G\"ottingen,
309--316 (1898).

\bibitem{BinetCauchy}
J.~Binet, J.~\'Ec.~Polytech. \textbf{9},
280 (1812);
A.~L.~Cauchy, ibid. \textbf{10}, 29 (1815).

\bibitem{Euler}
L.~Euler,
``De summis serierum reciprocarum,''
Comment.~Acad.~Sci.~Petropol. \textbf{7},
123--134 (1740).

\bibitem{Mehta}
M.~L.~Mehta,
\emph{Random Matrices}, 3rd ed.,
Academic Press (2004).

\bibitem{HardyWright}
G.~H.~Hardy and E.~M.~Wright,
\emph{An Introduction to the Theory of Numbers},
6th ed., Oxford (2008).

\end{thebibliography}

\bigskip
\noindent\textbf{Joint research by Mingu Jeong
and Claude (Anthropic).}

\end{document}
